\subsection{复形态射的扩张}

引理1. --- 考虑一个由A-模及同态构成的图表

\[
\begin{array}{c} \mathbf {M} ^ {\prime} \xrightarrow {\alpha^ {\prime}} \mathbf {M} \xrightarrow {\alpha} \mathbf {M} ^ {\prime \prime} \\ f ^ {\prime} \Biggl \downarrow \quad \swarrow k / f \Biggl \downarrow \quad \swarrow k ^ {\prime \prime} \\ \mathbf {N} ^ {\prime} \xrightarrow [ \beta^ {\prime} ]{} \mathbf {N} \xrightarrow [ \beta ]{} \mathbf {N} ^ {\prime \prime} \end{array}
\]

使得 \(f \circ \alpha' = \beta \circ f', \alpha \circ \alpha' = 0\) ，Ker \(\beta = \operatorname{Im} \beta'\) ，且 \(f = k'' \circ \alpha + \beta \circ k\) 且其中 \(\mathbf{M}'\) 是投射的。存在一个A-同态 \(k': \mathbf{M}' \to \mathbf{N}'\) 使得 \(f' = k \circ \alpha' + \beta' \circ k'\) .

事实上，令 \(g = f' - k \circ \alpha'\) ；我们有

\[
\beta \circ g = \beta \circ f ^ {\prime} - \beta \circ k \circ \alpha^ {\prime} = f \circ \alpha^ {\prime} - \beta \circ k \circ \alpha^ {\prime} = k ^ {\prime \prime} \circ \alpha \circ \alpha^ {\prime} = 0.
\]

这意味着 \(\operatorname{Im}(g) \subset \operatorname{Ker}(\beta) = \operatorname{Im}(\beta')\) 。由于 \(\mathbf{M}'\) 是投射的，因此存在一个A-同态 \(k': \mathbf{M}' \to \mathbf{N}'\) 使得 \(\beta' \circ k' = g\) ，从而引理得证。

引理2. --- 如果在A-模与同态的交换图中

\[
\begin{array}{c c c c c} \mathbf {M} ^ {\prime} & \xrightarrow {\alpha^ {\prime}} & \mathbf {M} & \xrightarrow {\alpha} & \mathbf {M} ^ {\prime \prime} \\ & & \Big \downarrow^ {u} & & \Big \downarrow^ {u ^ {\prime \prime}} \\ \mathbf {N} ^ {\prime} & \xrightarrow {\beta^ {\prime}} & \mathbf {N} & \xrightarrow {\beta} & \mathbf {N} ^ {\prime \prime} \end{array}
\]

我们有 \(\alpha\circ\alpha'\) =0，Ker \(\beta=\operatorname{Im}\beta'\) 且如果 \(\mathbf{M}'\) 是投射的，则存在一个A-同态 \(u':\mathbf{M}'\to\mathbf{N}'\) 使得 \(\beta'\circ u'=u\circ\alpha'\) .

只需在引理1中设 \(k'' = u'', k = -u, f = 0, f' = 0\) 与 \(u' = k'\) 即可。

引理1 bis. --- 考虑一个A-模与同态的图表

\[
\begin{array}{c} \mathbf {M} ^ {\prime} \xrightarrow {\alpha^ {\prime}} \mathbf {M} \xrightarrow {\alpha} \mathbf {M} ^ {\prime \prime} \\ \Big \downarrow_ {k ^ {\prime}} \Big \downarrow_ {f ^ {\prime}} \Big \downarrow_ {k} \Big \downarrow_ {f} \\ \mathbf {N} ^ {\prime} \xrightarrow [ \beta^ {\prime} ]{} \mathbf {N} \xrightarrow [ \beta ]{} \mathbf {N} ^ {\prime \prime} \end{array}
\]

使得 \(f \circ \alpha' = \beta \circ f'\) ，Ker \(\alpha = \operatorname{Im} \alpha'\) , \(\beta \circ \beta' = 0\) ，且 \(f' = k \circ \alpha' + \beta' \circ k'\) ，且其中 \(N''\) 是单射的。存在一个 A-同态 \(k'' : M'' \to N''\) 使得 \(f = k'' \circ \alpha + \beta \circ k\) 。事实上，设 \(q = f - \beta \circ k\) ，我们有

\[
g \circ \alpha^ {\prime} = f \circ \alpha^ {\prime} - \beta \circ k \circ \alpha^ {\prime} = \beta \circ f ^ {\prime} - \beta \circ k \circ \alpha^ {\prime} = \beta \circ \beta^ {\prime} \circ k ^ {\prime} = 0.
\]

这蕴含Ker \(g \supset \operatorname{Im} \alpha' = \operatorname{Ker} \alpha\) 。由于 \(\mathbf{N}''\) 是单射的，因此（X，第16页，注记）存在一个A-同态 \(k': \mathbf{M}'' \to \mathbf{N}''\) 使得 \(g = k'' \circ \alpha\) ，从而引理得证。

引理2 bis. --- 若在A-模与同态的交换图中

\[
\begin{array}{c c c c c} \mathbf {M} ^ {\prime} & \stackrel {{\alpha^ {\prime}}} {{\longrightarrow}} & \mathbf {M} & \stackrel {{\alpha}} {{\longrightarrow}} & \mathbf {M} ^ {\prime \prime} \\ u ^ {\prime} \Big \downarrow & & u \Big \downarrow & & \\ \mathbf {N} ^ {\prime} & \stackrel {{\beta^ {\prime}}} {{\longrightarrow}} & \mathbf {N} & \stackrel {{\beta}} {{\longrightarrow}} & \mathbf {N} ^ {\prime \prime}, \end{array}
\]

我们有 \(\operatorname{Ker} \alpha = \operatorname{Im} \alpha'\) , \(\beta \circ \beta' = 0\) 且如果 \(\mathbf{N}''\) 是单射的，则存在一个A-同态 \(u'' : \mathbf{M}'' \to \mathbf{N}''\) 使得 \(u'' \circ \alpha = \beta \circ u\) .

只需在引理 1 bis 中设 \(u' = k', u = -k, f = 0, f' = 0\) 与 \(k'' = u''\) 即可。

命题1. --- 设(P, \(d_{\mathbb{P}}\) ) 和 (E, \(d_{\mathbb{E}}\) ) 是 A-模的两个复形，且 r 是一个整数。

\begin{enumerate}
\def\labelenumi{\alph{enumi})}
\item
  设 \((u_i: \mathrm{P}_i \to \mathrm{E}_i)_{i \leqslant r}\) 是一个同态族，使得 \(d_{\mathrm{E}} \circ u_i = u_{i-1} \circ d_P\) （对 \(i \leqslant r\) ）。假设 \(\mathrm{P}_i\) 对 \(i > r\) 是投射的，且 \(\mathrm{H}_i(\mathrm{E}) = 0\) （对 \(i \geqslant r\) ）。则 \(u_i\) 的族可延拓为从 P 到 E 的一个复形态射；两个这样的延拓是同伦的。
\item
  设 \((u^{i}:\mathbf{P}^{i}\to\mathbf{E}^{i})_{i\leqslant r}\) 是一个同态族，使得 \(u^{i}\circ d_{\mathrm{P}}=d_{\mathrm{E}}\circ u^{i-1}\) （对 \(i\leqslant r\) ）。假设 \(\mathbf{E}^{i}\) 对 \(i>r\) 是单射的，且 \(\mathbf{H}^{i}(\mathbf{P})=0\) （对 \(i\geqslant r\) ）。则 \(u^{i}\) 的族可延拓为从 \(\mathbf{P}\) 到 \(\mathbf{E}\) 的一个复形态射；两个这样的延拓是同伦的。
\end{enumerate}

让我们证明 a)。族 \((u_{i})_{i\leqslant r}\) 的一个扩张 v 的存在性由引理 2 通过归纳法直接得出。设 \(v'\) 是另一个扩张；设 \(f=v'-v\) ，并对整数 n 作归纳，构造一个同态 \(k_{n}:P_{n}\to E_{n+1}\) ，使得 \(f_{n}=d_{E}\circ k_{n}+k_{n-1}\circ d_{P}\) 。对 \(i\leqslant r\) ，取 \(k_{i}=0\) 。设 \(n\geqslant r\) 并假设 \(k_{i}\) 对 \(i\leqslant n\) 已构造。于是考虑图表

\[
\begin{array}{c} \mathrm{P} _ {n + 1} \xrightarrow {d _ {\mathrm{p}}} \mathrm{P} _ {n} \xrightarrow {d _ {\mathrm{p}}} \mathrm{P} _ {n - 1} \\ f _ {n + 1} \Biggl \downarrow \quad k _ {n} \Biggl \downarrow f _ {n} \Biggl \downarrow k _ {n - 1} \\ \mathrm{E} _ {n + 2} \xrightarrow {d _ {\mathrm{E}}} \mathrm{E} _ {n + 1} \xrightarrow {d _ {\mathrm{E}}} \mathrm{E} _ {n}. \end{array}
\]

引理1的假设得到满足；因此存在一个A-同态 \(k_{n+1}: \mathbf{P}_{n+1} \to \mathbf{E}_{n+2}\) 使得 \(f_{n+1} = d_{\mathbf{E}} \circ k_{n+1} + k_n \circ d_p\) ，从而 \(a\) ).

b) 的证明类似，借助引理 1 bis 与引理 2 bis 完成。

